\begin{abstract}
テニスの単眼画像からコートを推定する課題は、白線や交点という局所的な画像証拠と、
その証拠を一つの射影幾何へまとめる大域的な整合性の両方を要求する。
キーポイントのヒートマップや線マスクは2次元では安価に教師を作れるが、
なぜその交点がその画素に落ちるのかは教えない。
そのため検出器は、局所的には妥当でも物理的には異なるコートへ対応する解を選びうる。
一方、画像ごとに正しいカメラ校正を付与することは、画像ごとに人手で幾何を測ることを意味する。

本稿はこの非対称性を、実会場の3Dガウススプラッティング（3DGS）再構成を使って移動させる。
提案する\Method では、
まず画像から抽出したコート線証拠を推定した地上面へ持ち上げ、
規制寸法のコートテンプレートを共有メートル尺度で当てはめ、
得られた剛体変換を学習に使わないカメラで一度だけ検証する。
この検証を通ったシーンだけが、シーン--コート変換の権威として公開される。
次に、復元カメラの凸包で解析的に制限した三次元軌道から新規視点を生成し、
同一の変換からキーポイント、領域、線、線の意味分類、カメラ姿勢を同時に作る。
人手の注釈費用は画像ごとではなく、シーンごとに一度のアライメント検証へ移る。

本研究の主たる問いは、
実会場の外観を継承しつつ撮影位置・方向・可視範囲の分布を広げた新規視点が、
放送映像中心の実データでは不足しがちな視点条件へテニスコート検出を適応させられるか、
である。
本稿はこの問いに対して、
視点分布を制御した再構成由来の教師生成、
単一のシーン--コートアライメントから複数教師を整合生成する契約、
および比較相手へも適用する共通のアライメント評価手続きを提示する。
画像数の増加、3DGS の利用、複数ヘッド化そのものを新規性とは主張しない。

本稿はこの設計を実測で裏づける。
公開実装ベースラインの生のキーポイントにも同じ RANSAC アライメントを適用し、
実写 held-out validation と合成 trajectory-group test の
それぞれで決定論的 \resBenchDomainSamples 件（seed \resBenchSeed）を
全 CPU で比較した。
実写 validation では \Method が僅かに上回り、
PCK@0.01 は \resRealOursPckOne 対 \resRealTcdPckOne、
対応誤差の中央値は \resRealOursPairMedian px 対 \resRealTcdPairMedian px である。
合成 test では差が大きく、
\Method の PCK@0.01 が \resSynthOursPckOne であるのに対し、
ベースラインは \resSynthTcdPckOne で、
有効なキーポイント対応は \resSynthTcdPairN 点にとどまった。
ただしこの差は、アーキテクチャと学習データが同時に異なる条件で測ったものであり、
3DGS による視点拡張の因果効果を分離しない。
実写のみで学習した同一アーキテクチャ、
同枚数の狭域・広域比較、未知会場の分離、複数 seed は未実施である。
checkpoint 選択時に記録された合成 validation の値は別結果として扱い、
交差モデル比較と混同しない。
\end{abstract}
